\section{Discussion}
\label{sec:discussion}

Using population-scale differentiable simulation, we have shown that
the reduced-order fiber sampling standard in GPU-accelerated
selectivity modeling --- reducing each fascicle to one or a few
representative fibers --- systematically overestimates the selectivity
actually delivered to the nerve.  The error is not a fixed offset: this
\emph{deployment penalty} is governed by fascicular architecture, the
number of fibers per target fascicle, and is therefore small in swine
(median $0.03$) but large in human (median $0.27$), a translational
bias that reduced-order models cannot detect because they never
simulate the fibers the bias acts on.  Mechanistically the penalty is
a failure to recruit the large human target fascicle from a field
tuned to a few sampled fibers, not off-target spill or reduced drive.
This was measurable only because jaxon simulates the full
${\sim}1000$-fiber population: a fully differentiable, GPU-native
reimplementation of the canonical MRG, Sweeney, Sundt and Rattay axon
models that reproduces NEURON ($99.6\,\%$ of $943$ configurations under
the $1\,\%$ criterion) at a geometric-mean $\sim\!820\times$ throughput
gain (peak $1\,377\times$) on a single A100, and exposes end-to-end
gradients over per-contact stimulus amplitudes.  Those gradients are
the enabling capability: previous PNS pipelines have been
gradient-free, restricting selectivity search to grid sweeps or
evolutionary methods, whereas jaxon recovers selective configurations
directly on subject-specific microCT vagus anatomies
(Section~\ref{sec:results-duke}) --- and, crucially, can re-score them
on the full population to expose the penalty
(Section~\ref{sec:results-sparse}).

\subsection{Comparison with prior work}

The closest direct comparator is pyfibers~\cite{hussain2024}, which
established the convention of wrapping NEURON axon models from Python
and using random-search / FD-gradient optimization for selectivity
design.  jaxon's contribution relative to pyfibers is twofold:
(i) the JAX implementation removes the Python-NEURON IPC overhead
that limits pyfibers throughput and makes vmap-batched parallel-fiber
forward passes possible on a single GPU; and (ii) the JAX implementation
admits end-to-end gradient computation, which pyfibers does not.
The validation results in
Section~\ref{sec:results-validation} confirm that this re-derivation
matches NEURON's accuracy on the canonical benchmark suite, so the
throughput and gradient gains do not come at the cost of model
fidelity.

The wider differentiable-simulation context for jaxon spans both
the JAX-ecosystem frameworks
\texttt{jaxley}~\cite{deistler2024} and
\texttt{JAX MD}~\cite{schoenholz2020jaxmd}, and the broader
autodiff-through-ODE literature initiated by the Neural ODE
work~\cite{chen2018neuralode}.  \texttt{jaxley} provides the
underlying multi-compartment HH machinery used by jaxon's solver.  \texttt{jaxley}
ships canonical CNS Hodgkin--Huxley channels and a generic cable
solver but does not include the calibrated peripheral nerve fiber
models that the PNS modeling literature relies on; jaxon's contribution
within this ecosystem is the peripheral specialization: the coupled
$(V_i, V_{\mathrm{px}})$ MRG double-cable solver, the calibrated
Sweeney, Sundt and Rattay fiber models, the validation-against-NEURON
harness, and the selectivity-optimization pipeline are domain-specific
and not provided out-of-the-box by \texttt{jaxley}.

GPU-accelerated multi-compartment neuron simulation has a separate
prior literature outside the autodiff context.
\texttt{NeuroGPU}~\cite{benshalom2022neurogpu} ports NEURON
biophysical models to CUDA and reports up to $200\times$ speedup over
NEURON-MPI on a single GPU, and
\texttt{Brian2GeNN}~\cite{stimberg2020brian2genn} accelerates Brian
spiking-network simulations via the GeNN code generator with
comparable throughput gains.  Crucially, however, none of these
GPU simulators models \emph{extracellular} field stimulation, which is
the regime relevant to cuff-electrode neuromodulation: \texttt{NeuroGPU}
supports intracellular current injection only,
\texttt{CoreNEURON}~\cite{kumbhar2019coreneuron} does not implement
NEURON's \texttt{extracellular} mechanism, and \texttt{Arbor} models
no extracellular fields --- so CPU-based NEURON has remained the only
established simulator of extracellular peripheral-fiber stimulation.
This is why our throughput comparison is against CPU NEURON, and why
jaxon's GPU-native extracellular forward pass is itself a capability
gap that the existing GPU simulators do not fill.  These pipelines
also accelerate only the forward pass and, like pyfibers, expose no
end-to-end gradients; gradient-based optimization over design
variables still requires the finite-difference detour the present work
is designed to obviate.

\subsection{Limitations}

\paragraph{Scope of the deployment-penalty finding.}
The penalty is established \emph{in silico}, at a single fiber diameter
($5.7\,\mu\mathrm{m}$ MRG), with one cuff and one FE pipeline, on 29
nerves (11 human).  We therefore frame fascicular architecture as the
\emph{driver} the penalty scales with (per-nerve Spearman
$r = +0.72$; fibers-per-fascicle mediates the species gap), not as a
manipulated cause, and we expect the magnitude --- though not the
direction --- to depend on diameter, fiber-type mix, cuff geometry and
the conductivity model.  Confirming the effect across diameters and
fiber populations, and anchoring it experimentally, is the natural next
step; jaxon's throughput makes the required population-scale sweeps
tractable.

\paragraph{Single-cuff geometry.}
The microCT selectivity results use a single 12-contact MultiContact
cuff (three axial rows $\times$ four angular columns) at a fixed
radius.  Larger contact counts, conformable cuffs and multi-cuff
arrays are increasingly common in clinical practice but were out of
scope for the present paper; jaxon's per-contact amplitude
parameterization extends to arbitrary contact layouts with no change
to the solver.

\paragraph{Simplified cuff finite-element model.}
The per-contact $V_e$ lead-fields for the microCT cohort are computed
with a simplified cylindrical-cuff finite-element model rather than a
fully device-faithful encapsulation/electrode-impedance solve.  jaxon's
interface accepts $V_e$ as a pre-computed
$[K, n_{\mathrm{fiber}}, n_{\mathrm{comp}}]$ tensor and so is agnostic
to its origin, so a higher-fidelity FEM can be substituted without
touching the optimization pipeline.

\paragraph{No biological variability in fiber parameters.}
All fibers in a given diameter share identical channel densities,
geometric ratios and resting state.  Inter-fiber variability is known
to broaden activation distributions and lower achievable
selectivity~\cite{aristovich2021}.  Future work will sample from
diameter-stratified parameter distributions estimated from
electrophysiological recordings.

\paragraph{Activation proxy versus full spike detection.}
The activation proxy
\begin{equation*}
  a_f = \min\!\bigl(\max_t \mathbf{m}_{\mathrm{Na}}(t)\big|_{\text{first node}},\
                    \max_t \mathbf{m}_{\mathrm{Na}}(t)\big|_{\text{last node}}\bigr)
\end{equation*}
correctly classifies recruitment when an action potential
propagates to both ends.  It is approximate in pathological regimes
(failed propagation, partial recruitment in dendrites with branching
or blocked nodes).  None of these regimes was encountered in the
straight-axon MRG simulations of this paper, but they would warrant
revisiting the proxy for branched-fiber or blocked-conduction
applications such as DC block or kHz conduction
block~\cite{kilgore2014block}.

\subsection{Outlook}

The combination of GPU throughput and end-to-end gradients positions
jaxon for several natural next-step applications:
(i) fiber-type (diameter-based) selectivity, exploiting the
differentiable forward model and waveform-shape gradients to recruit
one diameter population over another;
(ii) population-averaged designs that marginalise over realistic
nerve-to-nerve variability via a single batched forward pass;
(iii) inverse modeling of recorded extracellular field data to
infer fiber recruitment from cuff-recorded compound action
potentials~\cite{vissamsetti2025cap}; and
(iv) closed-loop adaptive stimulation policies trained by
backpropagation through the forward model.  Each of these would
require a gradient-free pipeline to perform thousands or millions
of redundant forward simulations; with jaxon, they reduce to standard
gradient-based optimization.
